\scanpage{3}
\reversemarginpar
\llap{\tikz[baseline=-.2ex]{\fill (0,0) circle (.018);}\hspace{1.3em}}12)\margin{commencé par Raynaud}% slanted
\hspace{.4em}Mise au point, et extensions diverses, de la théorie des modèles de Néron. \hfill
\begin{minipage}[c]{1.8em}\figure{\tikz[baseline=-.35ex,x=1ex,y=1ex]{\draw (.1,2) -- (0,-.5); \draw (.36,1.85) -- (.24,-.45);}}\end{minipage}

\llap{\tikz[baseline=-.2ex]{\draw[line width=.7pt] (0,.02)--(.045,-.015);}\hspace{1.3em}}\textcircled{\small 13}\margin{S'impose moins depuis les résultats d'approximation d'Artin}% slanted
\hspace{.4em}Théorémes \unsure{assymptotiques|asymptotiques} pour les foncteurs de Greenberg, relatifs à un anneau local noeth complet ou hensélien.

\llap{\tikz[baseline=-.2ex]{\draw[line width=.7pt] (0,.015)--(.06,0);}\hspace{1.3em}}14) Structures des Topos (questions d'existence de foncteurs fibres....). Structure des champs champs représentables.... \hfill
\begin{minipage}[c]{1.8em}\figure{\tikz[baseline=-.4ex,x=1ex,y=1ex]{\draw[double,double distance=.3pt] (.7,2) -- (1,2.1) -- (.46,-1) -- (.13,-.65);}}\end{minipage}

\llap{\tikz[baseline=-.2ex]{\draw[line width=.7pt] (0,0)--(.015,.045)--(.04,.04);}\hspace{1.3em}}\textcircled{\small 15}\margin{fait partiellement par Mazur-Artin. Réflexions de Quillen en train}% slanted
\hspace{.4em}\unsure{Théorie|Théories} du \deleted{\unsure{g}}type d'homotopie des Topos, incluant une définition des $\pi_{\unsure{i|\bullet|\cdot}}$ à \deleted{coefficients} noveaux infinitésimaux. \inserted{Th's d'\unsure{Hurwisz|Hurwicz}. Suites exactes d'homotopie. Ths de comparaison, \deleted{de \unsure{spécialisation}.}}

\llap{\tikz[baseline=-.2ex]{\draw[line width=.7pt] (0,.035)--(.01,0);}\hspace{1.3em}}\textcircled{\small 16}\margin{contre-exemple de Artin, situation peu engageante}% slanted
\hspace{.4em}\deleted{Théorie de spécialisation du} \tikz[baseline=(vr.base)]{\node[inner sep=1pt] (vr) {\inserted{``vrai''}};\draw ([xshift=-2pt]vr.south west) .. controls +(0,-.15) and +(-.2,0) .. ([yshift=-3pt]vr.south) -- ([yshift=-8pt]vr.south) -- ([xshift=2pt,yshift=-3pt]vr.south) .. controls +(.2,0) and +(0,-.15) .. ([xshift=2pt]vr.south east);} \deleted{groupe fondamental, pour une famille de courbes.} \hfill
\begin{minipage}[c]{1.8em}\figure{\tikz[baseline=-.4ex,x=1ex,y=1ex]{\draw (.75,1.65) -- (1,1.65) -- (.67,-.8) -- (.48,-.65);\draw (.93,.65) -- (.61,-.93) -- (.43,-.72);}}\end{minipage}

\llap{\tikz[baseline=1.1cm]{\draw (.13,1.25)--(0,1.2)--(.03,0)--(.18,.03);}\hspace{.25em}}17)\margin{résolu par Raynaud}% slanted
\hspace{.4em}\deleted{Séparation de foncteurs : classes de correspondance.... Cas de représentabilité et de non représentabilité de \mbox{$\underline{\mathrm{Pic}}_{X/S}$}, \mbox{$X$} de dim} \deleted{\unsure{de|rel}} \deleted{relative \mbox{$1$} sur \mbox{$S$}.}

\llap{\tikz[baseline=-.2ex]{\draw[line width=.7pt] (0,.04)--(.03,0);}\hspace{1.3em}}\textcircled{\small 18} Question du passage au quotient dans cas \unsure{arbitraire|artinien} et de l'effectivité de données de descente plates. \deleted{\unsure{Construction|Constructions} \unsure{générale|générales} de corps de \unsure{modules|variables}.} \hfill
\begin{minipage}[c]{1.8em}\figure{\tikz[baseline=-.4ex,x=1ex,y=1ex]{\draw[double,double distance=.25pt] (.9,1.8) .. controls (1,1.4) and (.7,0) .. (.31,-.8) -- (.08,-.5);}}\end{minipage}

\llap{\tikz[baseline=-.2ex]{\draw[gray!70,line width=.5pt] (0,0)--(.015,.04);}\hspace{1.3em}}\tikz[baseline=(num.base)]{\node[draw,circle,inner sep=.7pt] (num) {\shortstack{19\\[-.5ex]\scriptsize$\unsure{\times|*}$}};}
Théorie d'amplitude pour cycles de dim. quelconque.
% archive sheet number: 2
\normalmarginpar

